% Shared preamble for CS102 quizzes and exams.
% \examlabel (set by the wrapper) is the label in the page header, e.g. "Quiz 4" or "Midterm".
% \ifanswers controls whether \cb{1} shows as checked and \ifanswers{...} text
% renders. The build script's wrapper files set it explicitly with
% \answerstrue / \answersfalse after \input-ing this preamble.

\documentclass[11pt]{article}
\usepackage[margin=0.85in]{geometry}
\usepackage{fancyhdr}
\usepackage{amssymb}
\usepackage{enumitem}
\usepackage{booktabs}
\usepackage{array}
\usepackage{listings}
\usepackage{xcolor}
\usepackage{parskip}

\newif\ifanswers
\answerstrue

\pagestyle{fancy}
\setlength{\headheight}{14pt}
\fancyhf{}
\fancyhead[L]{CSCI-102 Data Structures - Fall 26 - \examlabel}
\fancyhead[R]{NetID: \underline{\hspace{1.6in}}}
\renewcommand{\headrulewidth}{0pt}
\fancyfoot[C]{\thepage}

\lstdefinestyle{java}{
  language=Java,
  basicstyle=\ttfamily\small,
  keywordstyle=\color{blue!55!black}\bfseries,
  commentstyle=\color{green!40!black}\itshape,
  stringstyle=\color{orange!70!black},
  showstringspaces=false,
  breaklines=true,
  backgroundcolor=\color{black!5},
  frame=none,
  xleftmargin=1em,
  framexleftmargin=1em,
  columns=fullflexible,
  keepspaces=true,
}
\lstset{style=java}

% \cb{1} = correct answer (renders checked only when answers are on), \cb{0} = incorrect
\newcommand{\cb}[1]{\ifanswers\ifnum#1=1 $\boxtimes$\else$\square$\fi\else$\square$\fi}
% \ansbox{k} = empty box to write a number in; shows k when answers are on
\newcommand{\ansbox}[1]{\framebox[2em][c]{\rule{0pt}{1.6em}\ifanswers #1\fi}}

\newcounter{qnum}
% separator rule between consecutive questions (none after the last one);
% drawn as glue (\cleaders) so it is discarded instead of starting a new page
\newcommand{\qsep}{\par\bigskip
  \cleaders\hbox to \linewidth{\hfil\rule{0.4\linewidth}{0.4pt}\hfil}\vskip 0.4pt
  \medskip}
\newenvironment{question}{%
  \par\noindent
  \ifnum\value{qnum}>0 \qsep\fi
  \stepcounter{qnum}%
  \begin{minipage}{\linewidth}
  \noindent\textbf{\theqnum.}\enskip\ignorespaces
}{%
  \end{minipage}\par
}

% fixed-height answer area, identical in both versions so the key keeps the student layout;
% put the key's answer inside it with \ifanswers ... \fi (it must fit in the given height)
\newenvironment{answerspace}[1]{\par\noindent\begin{minipage}[t][#1][t]{\linewidth}}{\end{minipage}}

\newcommand{\fullname}{Full Name: \underline{\hspace{4in}}\par\bigskip}
